\label{sec:validation}

With no wave buoy or other ground truth available for this dataset, validation relied on cross-checking the two independent wave estimators against each other, and on cross-checking the navigation/altitude pipeline against an independent reference signal. Each of the cases below follows the same pattern: a method was applied, a measurement revealed a quantitatively implausible or inconsistent result, a specific cause was diagnosed from the raw data, a fix was applied, and the corrected result was re-checked. Consistent with the characterize, implement and re-verify process described in the Methodology.

\textbf{Altimeter spectrum noise floor.} The first live run of the altimeter-based wave estimator produced a significant wave height of $\approx$2.3~m at a peak period of $\approx$1.7~s, a combination that implies a wave steepness well past the physical breaking limit, and therefore not a real sea state. Comparing the altimeter's elevation spectrum against the navigation filter's own altitude spectrum over the same window showed that the altimeter range carries a broadband sensor-noise floor that does not roll off above the real wave-frequency peak; integrating over the full available frequency band (rather than just the wave band) picked up that noise and inflated the height estimate roughly fourfold. Band-limiting the spectral integration to the physically plausible wave-frequency range corrected this to $H_s\approx0.4$--0.6~m at $T_p\approx2$--2.4~s, which is the range used throughout the rest of the validation. The same comparison also showed that the navigation filter's altitude signal itself carries essentially no energy in the wave band (the platform does not measurably heave at wave frequencies on this dataset), so subtracting it from the raw range, the originally intended heave-compensation step, was disabled, since it only injected the navigation filter's own noise without removing any real heave. Figure~\ref{fig:altimeter-pipeline} shows the full altimeter processing pipeline, with the band-limiting step highlighted.

\begin{figure}[H]
	\centering
	\includegraphics[width=0.72\textwidth,height=0.82\textheight,keepaspectratio]{diagrams/altimeter_wave_pipeline}
	\caption{Altimeter wave-estimation processing pipeline. The band-limiting step (highlighted) was the critical fix: integrating the full spectrum to Nyquist picked up the sensor noise floor and inflated $H_s$ roughly fourfold; restricting to the wave band [0.05--0.5~Hz] recovers the physical estimate.}
	\label{fig:altimeter-pipeline}
\end{figure}

\textbf{Point-cloud bad-frame contamination.} The point-cloud-based wave estimator initially produced an implausible maximum-wave-height estimate (2--9~m) at a significant height of only $\approx$0.5~m. Characterizing the surface-fit residuals frame by frame showed that roughly a third of frames suffered a whole-frame disparity failure (the entire reconstructed patch was spread out, not just a few stray points), so no single-frame outlier trimming could fix it. A per-frame quality gate was added that scores each frame's surface fit and discards the frame entirely if its implied wave height exceeds a physically reasonable threshold before it can contaminate the running estimate. This brought the height estimate down to a physically consistent value and improved the ratio between maximum and significant wave height from 4.6 (implausible) to 2.7 (consistent with real sea states). This per-frame quality gate is summarised in Figure~\ref{fig:bad-frame-gate}.

\begin{figure}[H]
	\centering
	\includegraphics[width=0.72\textwidth,height=0.82\textheight,keepaspectratio]{diagrams/bad_frame_gate}
	\caption{Per-frame quality gate in the point-cloud wave estimator. Whole-frame disparity failures (affecting $\approx$35\% of frames on the water surface) spread the entire reconstructed patch; the gate drops these before they contaminate the running $H_s$ estimate.}
	\label{fig:bad-frame-gate}
\end{figure}

\textbf{Reconciling the point-cloud height with the altimeter.} Even after the bad-frame fix, the point-cloud height estimate read roughly half the altimeter's. This was traced to the imaging geometry rather than a remaining bug: the reconstructed patch spans a long, oblique range, and far points carry disproportionately large stereo error, which both inflates a purely vertical variance estimate and biases a tilted-plane fit (which partly absorbs real long-wavelength elevation as ``tilt''). Restricting the analysis to the near field, and reporting the median rather than the mean over the buffered frames (the bad-frame gate otherwise biases the mean down asymmetrically), brought the point-cloud height estimate into agreement with the altimeter's at the same instant in the recording.

\textbf{Reconciling the period estimate via a Doppler/en\-counter-fre\-quency correction.} Even with height reconciled, the altimeter's peak period ($\approx$2.1--2.4~s) and the point cloud's spatially derived period ($\approx$3.7~s) disagreed by far more than measurement noise could explain. Recognizing that this dataset was recorded from a moving boat platform (at $\approx$2.6~m/s, up to 3.8~m/s) rather than a stationary one explained the gap: the altimeter measures the wave \emph{encounter} period as seen by a moving observer, not the true period, while the spatial point-cloud estimate is motion-independent. Applying the standard deep-water encounter-to-true frequency correction (using the platform's own speed, under a head-seas assumption) converted the altimeter's encounter period to a true period of $\approx$3.2~s, in agreement with the point cloud's independent 3.7~s estimate to within $\approx$10--13\%. This single platform-speed correction resolved what had been the longest-standing disagreement between the two estimators, and is consistent with the role of platform motion in altimetric wave measurement reported for other moving platforms~\cite{ichikawa2024seasurface}. Figure~\ref{fig:doppler-correction} illustrates this encounter-to-true period correction.

\begin{figure}[H]
	\centering
	\includegraphics[width=\textwidth]{diagrams/doppler_correction}
	\caption{Doppler/encounter-frequency correction. The altimeter, riding with the moving platform, measures the encounter period; the stereo point cloud measures the spatial (true) period directly. Applying the deep-water encounter-to-true correction reduces the period disagreement from $\approx$75\% to $\approx$10--13\%.}
	\label{fig:doppler-correction}
\end{figure}

\textbf{Wave direction.} A least-squares fit of cross-spectrum phase across the three altimeters (arranged roughly along the boat's longitudinal axis) was used to estimate an off-axis angle relative to head seas. The fit is internally consistent within a single analysis window (sub-degree residual) and the cross-spectrum coherence is high, confirming the array senses a genuinely coherent wave, but the inferred angle is small relative to the array's short aperture and drifted noticeably (3--20\textdegree) across different windows as the boat manoeuvred. The array therefore confirms a roughly head-seas geometry and is reported as a diagnostic quantity, but was not considered reliable enough to override the head-seas assumption used in the Doppler correction above, nor to resolve a full heading (a single line of altimeters cannot distinguish port from starboard, or fore from aft). Wave direction is accordingly not delivered as a pipeline output; this is a known shortfall relative to the original proposal's directional objective, and resolving it would require either a 2-D sensor array with sufficient aperture relative to the dominant wavelength, or a different sensing modality altogether (e.g.\ the 2-D FFT over the stereo point cloud, which was implemented but dropped due to 180° ambiguity and near-random instability on this patch size -- see Section~\ref{sec:future}).

\textbf{Overall assessment.} The two wave-parameter estimators are independent in both sensing modality (temporal, single-point altimetry vs.\ spatial, multi-point stereo geometry) and failure mode, and after the corrections above they agree on significant wave height ($\approx$0.4--0.6~m) and on peak period ($\approx$3.2~s vs.\ 3.7~s) to within $\approx$10--13\%. This mutual agreement, achieved without any external ground truth, is the strongest evidence in this report that the pipeline's wave-parameter outputs are physically meaningful rather than artifacts of either sensing modality.
